



\section{Experiments}
\subsection{Experimental Setting}
\label{ablation study}
\textbf{Datasets.}
We evaluate our method across two tasks: Multi-domain TIL (MTIL) and CIL. For MTIL, we follow the two-order training protocol proposed in \cite{zheng2023preventing}. For CIL, we follow~\cite{douillard2022dytox} to conduct experiments on CIFAR100 \cite{douillard2022dytox} and TinyImageNet \cite{yan2021dynamically}. The 100 classes of CIFAR100 are divided into $\{10, 20, 50\}$ subsets, and the 100 classes from TinyImageNet are divided into $\{ 5, 10, 20\}$ subsets to evaluate class distribution adaptability.

%


\noindent\textbf{Metrics.}
To evaluate our method on the MTIL, we utilize metrics proposed by \cite{zheng2023preventing}, namely ``Transfer'', ``Average'', and ``Last''. The ``Transfer'' metric assesses the model's zero-shot transfer capability on unseen data. ``Last'' evaluates the model's memorization ability on historical knowledge. ``Average'' is a composite metric measuring the mean performance across ``Transfer'' and ``Last''. In CIL, following \cite{douillard2022dytox}, we calculate the average accuracy over all subsets (``Average'') and specifically for the last subset (``Last'').

% 



\noindent\textbf{Implementation Details.}
As in \cite{zheng2023preventing}, we use the CLIP model with ViT-B/16~\cite{dosovitskiy2020image} as our backbone for all the experiments. 
We adopt LoRA~\cite{hu2021lora} as experts and set the total number $N_E=22$. The router is a single MLP that mixes the experts with $top$-$2$ gating scores. In DDAS, the reference data is TinyImageNet~\cite{yan2021dynamically}, and the threshold is $0.065$ and $0.06$ for full-shot and few-shot. The autoencoder is built upon a pretrained AlexNet~\cite{krizhevsky2012imagenet} with MLP and a non-linear layer. We use AdamW~\cite{loshchilov2017decoupled} optimizer and a label smoothing~\cite{muller2019does} technique for a better result.
For TIL, we train 1k iterations on full-shot and 500 iterations on few-shot for each task. For DDAS, we train 1k and 300 iterations for reference datasets and incremental tasks, respectively. Except for the reference dataset, the MoE-Adapters and DDAS are jointly trained during continual learning. %





\subsection{Comparison with State-of-the-art Methods}
\label{comparison}
\textbf{Multi-domain Task Incremental Learning.} 
Table~\ref{tab:compareAvg_Last_Transfer} showcases a comparison between our proposed method and alternative approaches in the MTIL task. Our approach utilizes the predefined Order-\uppercase\expandafter{\romannumeral1} from~\cite{zheng2023preventing}, where datasets are trained and tested sequentially in a left-to-right order, as displayed in Table~\ref{tab:compareAvg_Last_Transfer}. Additional results for Order-\uppercase\expandafter{\romannumeral2} are provided in the supplementary material. 
%
The uppermost section of Table~\ref{tab:compareAvg_Last_Transfer} displays the outcomes of applying CLIP independently on each task through zero-shot inference, full parameter fine-tuning, and parameter-efficient fine-tuning. The ``zero-shot" represents the optimal results of CLIP's zero-shot transfer on each task, while the other two rows indicate the highest possible outcomes achieved by fine-tuning CLIP in each respective task.
%
Our proposed method, labeled as ``Ours", outperforms the second-best approach on most tasks, resulting in an overall improvement of by 0.8\%, 1.3\%, and 1.4\% in terms of ``Transfer'', ``Average'' and ``Last'', respectively. Additionally, by increasing the training iterations from 1k to 3k, our method (labelled as ``Ours$\dagger$'') achieves a further improvement of 1.8\% and 3.8\% on ``Average'' and ``Last'' while dropping 0.7\% on ``Transfer'' in comparison to ZSCL~\cite{zheng2023preventing}. Furthermore, our model achieves less degradation than ZSCL when compared with the upper bound methods, demonstrating favorable performance in anti-forgetting and zero-shot transfer.  


\begin{table}[t]
        \centering
        \resizebox{0.95\linewidth}{!}{
	
	\begin{tabular}{lcc|cc|cc}
            \toprule 
            \multirow{2}{*}{Method} & \multicolumn{2}{c}{10 step} & \multicolumn{2}{c}{20 step} & \multicolumn{2}{c}{50 step} \\
            \cline{2-7} 
            \noalign{\smallskip}
			 & 	Avg. & Last &Avg. & Last &Avg. &Last  \\
		\midrule
            UCIR \cite{hou2019learning} & 58.66 & 43.39& 58.17 & 40.63 & 56.86 & 37.09 \\
            Bic\cite{wu2019large}&68.80 &53.54 &66.48&47.02 & 62.09 & 41.04\\
		PODNet\cite{douillard2020podnet} &58.03&41.05&53.97&35.02&51.19&32.99\\
            DER \cite{yan2021dynamically} & 74.64 & 64.35 & 73.98 & 62.55 & 72.05 & 59.76\\
            DyTox+\cite{douillard2022dytox} & 74.10 & 62.34 & 71.62 & 57.43 & 68.90 & 51.09 \\
            DNE \cite{hu2023dense} & 74.86& 70.04& - & - & -& - \\
            \midrule
            CLIP Zero-shot &74.47  & 65.92 & 75.20 & 65.74  & 75.67 & 65.94 \\
            Fine-tune & 65.46& 53.23& 59.69& 43.13&39.23 &18.89\\
            LwF \cite{li2017learning}&65.86 &48.04 &60.64 &40.56 &47.69 &32.90\\
            iCaRL  \cite{rebuffi2017icarl} & 79.35 &70.97&73.32 &64.55 & 71.28&59.07\\
            LwF-VR \cite{ding2022don}&78.81 &70.75 &74.54 &63.54 &71.02 & 59.45\\
            ZSCL \cite{zheng2023preventing} &\underline{82.15} &\underline{73.65} &\underline{80.39} &\underline{69.58} &\underline{79.92} &\underline{67.36}\\
        \rowcolor{Thistle!20}
		{Ours} & \textbf{85.21} &\textbf{77.52}& \textbf{83.72}&\textbf{76.20} & \textbf{83.60} & \textbf{75.24}  \\
		\bottomrule
	\end{tabular}}
    \vspace{-5pt}
	\caption{Comparison of different methods on CIFAR100 in class-incremental setting. We label the best and second-best methods with \textbf{bold} and \underline{underline} styles.}
	\label{tab:CIL1}
 \vspace{-9pt}
\end{table}
\begin{table}[t]
	\centering
         \resizebox{0.95\linewidth}{!}{
	\begin{tabular}{lcc|cc|cc}
            \toprule \multirow{2}{*}{Method} & \multicolumn{2}{c}{5 step} & \multicolumn{2}{c}{10 step} & \multicolumn{2}{c}{20 step} \\
            \cline{2-7} 
            \noalign{\smallskip}
		 & 	Avg. & Last &Avg. & Last & Avg. &Last  \\
		\midrule
            EWC \cite{kirkpatrick2017overcoming} & 19.01 & 6.00&15.82 &3.79 & 12.35&4.73  \\
            EEIL \cite{castro2018end} &47.17 & 35.12&45.03 &34.64 & 40.41&29.72\\
            UCIR \cite{hou2019learning}&50.30 & 39.42& 48.58& 37.29& 42.84& 30.85  \\
		  MUC \cite{liu2020more} &32.23&19.20 & 26.67& 15.33& 21.89&10.32  \\
            PASS \cite{zhu2021prototype} &49.54 & 41.64&47.19& 39.27& 42.01& 32.93 \\
            DyTox \cite{douillard2022dytox} & 55.58& 47.23&52.26 & 42.79& 46.18&36.21   \\
            
            \midrule
            CLIP Zero-shot & 69.62& 65.30& 69.55& 65.59& 69.49&65.30  \\
            Fine-tune &61.54 &46.66 &57.05 & 41.54& 54.62 & 44.55  \\
            LwF \cite{li2017learning}&60.97 &48.77 &57.60 &44.00 &54.79 &42.26  \\
            iCaRL  \cite{rebuffi2017icarl} & 77.02& 70.39& 73.48&65.97 &69.65 &64.68  \\
            LwF-VR \cite{ding2022don}& 77.56&70.89 &74.12 &67.05 & 69.94&63.89  \\
            ZSCL \cite{zheng2023preventing}& \underline{80.27}& \underline{73.57}&\underline{78.61} & \underline{71.62}& \underline{77.18}&\underline{68.30}  \\
        \rowcolor{Thistle!20}
		Ours& \textbf{81.12} & \textbf{76.81}& \textbf{80.23}& \textbf{76.35}& \textbf{79.96} & \textbf{75.77}   \\
		\bottomrule
	\end{tabular}}
 \vspace{-4pt}
	\caption{Comparison of different methods on TinyImageNet dataset in class-incremental settings with 100 base classes. We label the best and second methods with \textbf{bold} and \underline{underline} styles.}
	\label{tab:CIL2}
 \vspace{-11pt}
\end{table}



\noindent\textbf{Few-shot Multi-Domain Task Incremental Learning.} 
we also ran experiments on few-shot MTIL, limiting the CLIP model to access only a few samples per task. In a 5-shot setting, the comparison results are shown in Table~\ref{tab:comparefs} using the same metrics as Table~\ref{tab:compareAvg_Last_Transfer}. Our method outperforms most state-of-the-art approaches on most datasets, surpassing the second-best method by 3.6\%, 7.0\%, and 4.2\% in terms of ``Transfer," ``Average," and ``Last". These results demonstrate the effectiveness of the proposed incremental MoE-Adapters in addressing the forgetting issue in long-term continual learning, even with limited samples. Furthermore, our proposed DDAS effectively learns data distribution discrimination with fewer training samples.

\noindent\textbf{Class Incremental Learning.}
We conduct experiments on class incremental learning to verify our method's performance on single-domain CL. Unlike MTIL, the task id of the input image is unknown in CL. To this end, our MoE-Adapters use only one router with two experts to adapt to all the subsets. The comparison results between our method and state-of-the-art approaches on CIFAR100 and TinyImageNet are shown in Table~\ref{tab:CIL1} and~\ref{tab:CIL2}, respectively. As we can see, the proposed method consistently outperforms the other competitors, including dynamic expansion and CLIP-based approaches, demonstrating the effectiveness and scalability of our MoE-Adapters in addressing single-domain CL. 

\begin{table}[t]
	\centering
         \resizebox{0.96\linewidth}{!}{
	\begin{tabular}{c >{\centering\arraybackslash}p{2.4cm} >{\centering\arraybackslash}p{2.1cm}>{\centering\arraybackslash}p{1.8cm}}
		\toprule		Method & Train Params $\downarrow$ &GPU $\downarrow$ & Times $\downarrow$ \\
		\midrule
            \textcolor{Black}{LWF}~\cite{li2017learning}&\textcolor{Black}{149.6M} & \textcolor{Black}{32172MiB} & \textcolor{Black}{1.54s/it}\\
            \textcolor{Black}{LWF-VR}~\cite{ding2022don} & \textcolor{Black}{149.6M} & \textcolor{Black}{32236MiB} &\textcolor{Black}{1.51s/it}\\
            \rowcolor{gray!20}
		ZSCL~\cite{zheng2023preventing} &149.6M & 26290MiB & 3.94s/it \\
		MoE-Adapters &\ \ 51.1M & 19898MiB & 1.37s/it \\
		DDAS  & \ \ 8.7M & \ \ 2461MiB & 0.21s/it \\
            \rowcolor{Thistle!20}		
            Ours&	\ \ 59.8M&	22358MiB	& 1.58s/it \\
            \textbf{$\Delta$} &	\textbf{\textcolor{Maroon}{-60.03\%}}&	\textbf{\textcolor{Maroon}{-14.95\%}}	& \textbf{\textcolor{Maroon}{-59.90\%}} \\
		\bottomrule
	\end{tabular}}
	\caption{\textcolor{Black}{Comparison of computational cost during training between our method and others in terms of training parameters, GPU burdens and training times of each iteration. And the \textbf{$\Delta$} is the improvement relative to the SOTA ZSCL~\cite{zheng2023preventing}.}}
	\label{tab:cost}
 \vspace{-10pt}
\end{table}

\noindent\textbf{Computational Cost.}
The experiments above have demonstrated the promising performance of our method in both MTIL and CIL. We further compare the computational cost of our method with \textcolor{Black}{others} to prove its parameter and time efficiency during training. Table~\ref{tab:cost} shows that our method is superior to \textcolor{Black}{the SOTA method ZSCL}, with a reduction of approximately 60\%, 15\%, and 60\% in terms of training parameter (M), GPU burdens (MiB), and iteration time, respectively. Additionally, we analyze the efficiency of our two proposed components, MoE-Adapters and DDAS, demonstrating that they effectively enhance continual learning for CLIP while reducing significant computation burdens during training.  





% -------------------------------------------------------------------------------


\subsection{Ablation Study}
\label{ablation}
In this section, we mainly analyze the efficacy of the proposed incremental MoE-Adapters and DDAS. All the experiments are conducted in MTIL setting. More analysis can be found in the supplementary material.   

\noindent\textbf{Analysis of MoE-Adapters.}
We conduct detailed ablation studies of different settings on MoE-Adapters, as shown in Table \ref{tab:ablation1}. Compared with the zero-shot CLIP and the fine-tuned version with one adapter, our MoE-Adapters effectively mitigate the ``catastrophic forgetting'' issue and retain the zero-shot transfer ability on the unseen date. In the proposed MoE-Adapters, we use $T$ task-specific routers to adaptively activate the $Topk$ experts from the predefined expert pool. Table~\ref{tab:ablation1} illustrates several different experts and routers combinations. As we can see, compared with using more experts, the task-specific routers contribute more to improve anti-forgetting and zero-shot transfer abilities. In the training of MoE, we propose an incremental activate-freeze strategy, enabling the collaboration of previously learned experts and inactivated ones for more accurate prediction. The comparison between ``Ours'' and ``+22E/11R w/o F'' demonstrates the effectiveness of this strategy.  

%

\begin{table}[t]
	\centering
        \resizebox{0.95\linewidth}{!}{
	\begin{tabular}{lcc|cc|cc}
		\toprule		
            Method &  Transfer & $\Delta$ &Avg.&  $\Delta$ &  Last &$\Delta$ \\
		\midrule
		% CLIP {\footnotesize ViT-B/16@224px} & &  &  & &  &\\
            CLIP Zero-shot &  69.4 & \textcolor{Maroon}{+0.5} & 65.3 &  \textcolor{NavyBlue}{-11.4} &  65.3 &  \textcolor{NavyBlue}{-19.7}\\
		{\hspace{1em}}+Adapter &45.0 & \textcolor{NavyBlue}{-23.9}  & 57.0 & \textcolor{NavyBlue}{-19.7} &  71.5 & \textcolor{NavyBlue}{-13.5}  \\
		{\hspace{1em}}+2E/1R &45.1  &\textcolor{NavyBlue}{-23.8} &56.3 &\textcolor{NavyBlue}{-20.4}  & 71.1& \textcolor{NavyBlue}{-13.9}\\
            {\hspace{1em}}+2E/11R & 68.1 & \textcolor{NavyBlue}{-0.8} & 72.6&  \textcolor{NavyBlue}{-4.1}  & 77.9&  \textcolor{NavyBlue}{-7.1}\\
            {\hspace{1em}}+22E/1R & 44.1 & \textcolor{NavyBlue}{-24.8}&56.0 & \textcolor{NavyBlue}{-20.7} & 66.2& \textcolor{NavyBlue}{-18.8}\\
            {\hspace{1em}}+22E/11R w/o F & 68.6  &\textcolor{NavyBlue}{-0.3} & 75.1 &  \textcolor{NavyBlue}{-1.6} & 82.0 & \textcolor{NavyBlue}{-3.0}\\
		% \rowcolor{Thistle!20}
            % {\hspace{1em}}Ours$\dagger$ & 67.4& \textcolor{NavyBlue}{-2.0} & 77.2&  \textcolor{Maroon}{+11.9} & 87.4 & \textcolor{Maroon}{+22.1} \\
            \rowcolor{Thistle!20}
            {\hspace{1em}}Ours & 68.9& \ 0.0 & 76.7&  \ 0.0 &  85.0 & \ 0.0 \\
		\bottomrule
	\end{tabular}}
 \vspace{-5pt}
	\caption{Ablation studies on incremental MoE-Adapters. ``mE/nR'' indicates MoE with m experts and n routers, respectively. ``F'' represents the incremental activate-freeze strategy. }
  \vspace{-5pt}
	\label{tab:ablation1}
\end{table}

\begin{figure}[t]
	\centering
	\includegraphics[width=0.95\linewidth]{figs/figure_ablation.pdf}
 \vspace{-5pt}
	\caption{Analysis of expert's number in different training iterations. 
    The results can be referred to ``Ours'' and ``Ours$\dagger$'' in Table~\ref{tab:compareAvg_Last_Transfer}.
    }
	\label{fig:ablation}
  \vspace{-12pt}
\end{figure}

\noindent\textbf{Analysis of Expert Number.}
Figure~\ref{fig:ablation} presents the ablation study on the number of experts used in MoE-Adapters. The experiments are conducted on the models trained by 1k~/~3k iterations with $T$ task-specific routers. The smoothness of the curves in the figure indicates our method's robustness for changing the number of experts. We can observe that the three metrics remain relatively stable as the number of experts changes. 
As shown in Figure~\ref{fig:ablation} (a) and (b), it is more stable in the 1k iteration setting than in 3k iterations. This phenomenon arises due to the escalating number of iterations, leading to an increase in the overall frequency of expert selection. When the number of experts is small, our activate-freeze strategy may activate more untapped experts and train them a few times, leading to the router mistakenly activating incompletely trained experts during the inference phase. The fluctuation is within an acceptable range and does not significantly affect the final performance.

%\



\begin{figure}[t]
	\centering
	\includegraphics[width=1\linewidth]{figs/figure_tSNE.pdf}
 \vspace{-20pt}
	\caption{t-SNE on DDAS's output of each task on full-shot and few-shot MTIL. The corresponding task names from $id=1-11$ are matches with the datasets listed from left to right in Table~\ref{tab:compareAvg_Last_Transfer}. }
	\label{fig:tsne}
  \vspace{-12pt}
\end{figure}

\noindent\textbf{Analysis of DDAS.}
We use DDAS to automatically distinguish the input images from seen or unseen data by learning the variations in data distribution with task-specific autoencoders. To verify the effectiveness of DDAS, we analyze the distribution discrimination in feature space, whose results are illustrated in Figure~\ref{fig:tsne}. We employ the reconstructed features $\textit{\textbf{f}}^t_o$ and plot the figure when the continual learning is finished. As we can see, the proposed DDAS is effective at learning the discriminative distribution of each learned task in full-shot and few-shot MTIL. 
As shown in Figure~\ref{fig:tsne}, the feature distributions of some samples from Task 9 overlap with that of Task 11. It is because these samples are misclassified by DDAS as out-of-distribution data and perform feature extraction with reference autoencoder. Although the inevitable misclassifications occur, our method still outperforms the state-of-the-art approaches in various metrics.
%